\documentclass[12pt]{article}
\usepackage{amsmath,amssymb,amsthm,bm,geometry,hyperref,booktabs}
\geometry{margin=23mm}
\setlength{\parskip}{0.5em}
\setlength{\emergencystretch}{2em}
\newcommand{\ie}{\begin{equation}\begin{aligned}}
\newcommand{\fe}{\end{aligned}\end{equation}}
\newcommand{\NS}{\mathrm{NS}}
\newcommand{\R}{\mathrm{R}}
\numberwithin{equation}{section}
\title{A fermion insertion on the Ramond tube of the theta block}
\author{Separate working notes}
\date{September 10, 2026}
\begin{document}
\maketitle

We derive a fermion-inserted double-Virasoro representation of the
NS--R--R theta block with $\eta\eta'=-1$, using the states, BPZ pairing,
branching coefficients and spin signs of \texttt{Human Notes/SCblock.tex}.
Both Virasoro factors in the enlarged block are evaluated by
central-charge recursion.  The auxiliary factor is evaluated directly
by fermion-state sewing; its Ising decomposition identifies the
irreducible representations being propagated.
One contour phase missing from the current definition of $\hat\rho$ is
made explicit below.  It is necessary for the branching and physical
sewing conventions to give the stated convolution.

\section{The additional sphere and the operator on the tube}

The two original spheres have Neveu--Schwarz punctures at infinity and
Ramond punctures at one and zero.  We replace the second edge by a sphere
with coordinate $z_3$, with sewing equations
\ie
 (z_1-1)z_3=\hat q_{2,1},\qquad
 (z_2-1)z_3^{-1}=\hat q_{2,2}.
\fe
Eliminating $z_3$ gives
\ie
 (z_1-1)(z_2-1)=q_2,\qquad
 q_2=\hat q_{2,1}\hat q_{2,2}.
 \label{product}
\fe
Thus this insertion introduces a marked point on the second tube.  Moving
it changes the two sewing parameters while preserving their product.
We choose the transported Ramond spin frame on the additional sphere
so that the lifts satisfy $\eta_2=\hat\eta_{2,1}\hat\eta_{2,2}$.
This fixes the square-root transport convention in addition to the
coordinate relation \eqref{product}.

In the physical irreducible NS vacuum module,
$G_{-1/2}\mathbf1=0$.  The normalization in the human notes consequently
gives
\ie
 v_{1/2}(P=Q/2)=Q\psi_{-1/2}\mathbf1,
 \qquad V[v_{1/2}(Q/2)](1)=Q\psi(1),
 \qquad Q=b+b^{-1}.
\fe
The field has the physical identity as its SCA factor.  The insertion we
use is the radially ordered tube operator
\ie
 \Theta V[v_{1/2}(Q/2)](1)=Q\Theta\psi(1).
 \label{ordered-insertion}
\fe
Here $\Theta$ acts on the Ramond module on the outgoing side of the field;
it is not a prescription to apply a Ramond zero mode to an NS state.

For a chain $\boldsymbol\Psi_{-\mathsf B}$ of nonzero Ramond fermion
modes, define the odd Ramond intertwiner by
\ie
 \Theta(\boldsymbol\Psi_{-\mathsf B}u^{\mathsf b})
 =(-1)^{\#\mathsf B+\mathsf b+1}
 \boldsymbol\Psi_{-\mathsf B}u^{1-\mathsf b}.
\fe
The auxiliary conventions are
$\psi_0u^{\mathsf b}=u^{1-\mathsf b}/\sqrt2$,
$\langle u^{\mathsf b}|u^{\mathsf b}\rangle=(-1)^{\mathsf b}$ and
$\psi_r^\dagger=-\psi_{-r}$.  Directly from the displayed definition,
\ie
 \Theta^2=-1,\qquad \{\Theta,\psi_r\}=0,
 \qquad \{\Theta,G_r\}=0,\qquad
 [\Theta,L_m]=[\Theta,L_m^{(1)}]=[\Theta,L_m^{(2)}]=0.
\fe
The anticommutator with $G_r$ uses the graded tensor product.  The product
$\Theta\psi$ is even and acts only on the auxiliary factor.  Its
zero-mode part acts diagonally:
\ie
 Q\Theta\psi_0
 (\boldsymbol\Psi_{-\mathsf B}u^{\mathsf b})
 =\frac{Q}{\sqrt2}(-1)^{\mathsf b}
 \boldsymbol\Psi_{-\mathsf B}u^{\mathsf b}.
 \label{zero-mode-action}
\fe
The signs from moving $\psi_0$ past the nonzero modes and from $\Theta$
cancel.  This ground-state sign makes recovery of the $\eta\eta'=-1$
sector possible; the sector identity and inverse are derived below.
The uninserted auxiliary block annihilates that sector under convolution.
The full field also contains nonzero modes: we must not replace
$\Theta\psi(1)$ by $\Theta\psi_0$ inside a split-tube block.

\section{The double-Virasoro decomposition}

The physical parameters and central charge are
\ie
 h_1=\frac{Q^2/4-P_1^2}{2},\qquad
 \beta_2=\frac{P_2}{\sqrt2},\qquad
 \beta_3=\frac{P_3}{\sqrt2},\qquad c=\frac32+3Q^2.
\fe
Here $p_1$ is the intrinsic parity of the NS primary and $f$ is the
relative outer three-point label used in the human notes.  In particular,
the allowed descendant parities satisfy the selection rule displayed
in \eqref{full-branching}.  We assume $b\ne0$, $b^2\ne1$ for the
embedding and $Q\ne0$ for the chosen insertion and its sector inverse.
At isolated momentum poles the complete expressions are
understood by meromorphic continuation, rather than by separately
evaluating singular summands.

The internal double-Virasoro weights are
\ie
 h_{j,n}^{(1)}&=\frac{Q^2/4-(P_j+2nb)^2}{2(1-b^2)},&
 c^{(1)}&=1+\frac{3Q^2}{1-b^2},\\
 h_{j,n}^{(2)}&=\frac{Q^2/4-(P_j+2n/b)^2}{2(1-b^{-2})},&
 c^{(2)}&=1+\frac{3Q^2}{1-b^{-2}}.
\fe
The NS labels lie in $\frac12\mathbb Z$ and contribute level $2n^2$.
The Ramond labels lie in $\frac12\mathbb Z+\frac14$ and contribute level
$2n^2-\frac18$.  The latter subtraction measures level relative to the
Ramond ground states, including the auxiliary ground weight $1/16$.
The two weights of the inserted field are
\ie
 h_\psi^{(1)}=-\frac{1+2b^2}{2(1-b^2)},\qquad
 h_\psi^{(2)}=\frac{b^2+2}{2(1-b^2)}.
\fe
Their sum is $1/2$.  Its two level-two degenerate fusion rules require
\ie
 n'=n+\frac12\quad\hbox{or}\quad n'=n-\frac12.
 \label{middle-fusion}
\fe
The odd field reverses the Ramond parity and $\Theta$ reverses it back.
The two parts of the split edge therefore carry the same parity $\alpha$.
The outer three-point forms retain the same label $f$.

\subsection{The outer vertex convention}

The physical forms $\rho_f^{(\eta)}$ use the Ramond BPZ convention of
the human notes.  The free-field branching forms require the following
contour phase when they are expressed in that physical tensor basis:
\begin{align}
 &\hat\rho_f^{(\eta)}(
 \Psi_{-\mathsf A}\mathbf L_{-A}\phi_1,
 \boldsymbol\Psi_{-\mathsf B}u^{\mathsf b}\otimes\mathbb L_{-B}w_2^\alpha,
 \boldsymbol\Psi_{-\mathsf C}u^{\mathsf c}\otimes\mathbb L_{-C}w_3^\gamma)
 \nonumber\\
 &\quad=(-i)^{A\bmod2}
 (-1)^{f(A+B+|\alpha|)+A\mathsf A+(B+|\alpha|+p_1)(\mathsf C+\mathsf c)}
 \nonumber\\
 &\qquad\quad\times
 \rho_{\mathsf F}(\Psi_{-\mathsf A}\mathbf1,
 \boldsymbol\Psi_{-\mathsf B}u^{\mathsf b},
 \boldsymbol\Psi_{-\mathsf C}u^{\mathsf c})
 \rho_f^{(\eta)}(\mathbf L_{-A}\phi_1,
 \mathbb L_{-B}w_2^\alpha,\mathbb L_{-C}w_3^\gamma).
 \label{outer-frame}
\end{align}
In \eqref{outer-frame}, $|\alpha|$ is the parity of the physical Ramond
ground $w_2^\alpha$; in \eqref{full-branching},
$\alpha,\gamma\in\{0,1\}$ label the total parities of the enlarged
Ramond primaries.  Descendant indices in sign exponents count fermionic
generators; their absolute values denote levels.  The factor
$(-i)^{A\bmod2}(-1)^{f(A+B+|\alpha|)}$ follows by comparing the
NS--R--R supercurrent Ward identities in the free-field and physical
frames, after the Ramond ground basis is transported to its physical
BPZ normalization.  The implementation's three-point label must also be
transported as $\eta_{\mathrm{native}}=(-1)^f\eta$.
For example, in the odd form the relative coefficients of $\beta_2$
and $\beta_3$ in $\rho(G_{-1/2}\phi,w^+,w^+)$ have the wrong relative
sign if this label transport is omitted.
The NS supercurrent step changes the contour factor by
$-i(-1)^{A+f}$, while flipping the second Ramond parity changes it by
$(-1)^f$; substitution of these ratios reproduces the three physical
supercurrent Ward recursions.  The transport also has a direct
double-Virasoro interpretation: in this frame the physical supercurrent
and auxiliary fermion contours at infinity have multipliers $-i$ and
$+i$, so their product is $+1$, as required for the integer-spin field
$:G\psi:$ entering $L^{(1,2)}$.

Equation \eqref{outer-frame} supplies a correction to the current
human-notes factorization of $\hat\rho$, which omits this contour phase
while retaining $(-1)^{A+\mathsf A}$ in the enlarged sewing sum.
Without the phase, an unwanted $(-1)^A$ remains in the physical factor.
We use \eqref{outer-frame} throughout, including in the branching calculation.

\subsection{Sewing the four internal edges}

We first define the ordinary punctured Virasoro block, keeping the order
of the three punctures $(\infty,1,0)$ explicit.  In the following formula
the four internal weights are $(h_1,h_L,h_R,h_3)$; the additional weight
$h_\psi$ is external.  This use of $L,R$ only specifies which half of the
second tube is meant.  With each inverse Gram matrix restricted to equal
levels,
\ie
 &F(h_1,h_L,h_R,h_3;h_\psi;c\mid
 q_1,\hat q_{2,1},\hat q_{2,2},q_3)\\
 &\quad=\sum_{\substack{|A|=|A'|,\ |B|=|B'|\\
                         |C|=|C'|,\ |D|=|D'|}}
 q_1^{|A|}\hat q_{2,1}^{|B|}\hat q_{2,2}^{|C|}q_3^{|D|}
 G_{h_1}^{AA'}G_{h_L}^{BB'}G_{h_R}^{CC'}G_{h_3}^{DD'}\\
 &\qquad\quad\times
 \rho(L_{-A}\nu_{h_1},L_{-B}\nu_{h_L},L_{-D}\nu_{h_3})
 \rho(L_{-A'}\nu_{h_1},L_{-C}\nu_{h_R},L_{-D'}\nu_{h_3})\\
 &\qquad\quad\times
 \rho(L_{-C'}\nu_{h_R},\nu_{h_\psi},L_{-B'}\nu_{h_L}).
 \label{punctured-vir}
\fe
The leading coefficient is one.  There is no constraint $|B|=|C|$:
the insertion can change the auxiliary level.  The displayed Gram
contraction defines the block; the production computation evaluates it
by $c$-recursion, not by constructing these Gram matrices.

The middle primary matrix element needed for sewing is
\ie
 \langle v_{2,n'}^\alpha|\Theta V[v_{1/2}(Q/2)](1)
 |v_{2,n}^\alpha\rangle.
\fe
Commutation of $\Theta$ with the two Virasoro algebras implies that
every middle descendant matrix element is this primary matrix element
times the two ordinary forms in the last line of
\eqref{punctured-vir}.  The analogous statement holds at each outer
vertex.  Inserting the double-Virasoro decomposition on all four cuts
therefore gives the full formula
\begin{align}
 &\widehat{\mathbb F}_f^{(\eta,\eta')}[\Theta v_{1/2}]
 =\sum_{\substack{n_1\in\frac12\mathbb Z,\ n,n',n_3\in\frac12\mathbb Z+\frac14\\
 n'=n\pm\frac12,\ \alpha,\gamma\in\{0,1\}\\
 2n_1+\alpha+\gamma\equiv f\ (2)}}
 q_1^{2n_1^2}\hat q_{2,1}^{2n^2-1/8}
 \hat q_{2,2}^{2n'^2-1/8}q_3^{2n_3^2-1/8}
 \eta_1^{2n_1+p_1}\eta_2^\alpha\eta_3^\gamma\nonumber\\
 &\quad\times(-1)^{2n_1+(2n_1+p_1)(\alpha+\gamma)+\alpha\gamma}
 \frac{\hat\rho_f^{(\eta)}(v_{1,n_1}^{p_1},v_{2,n}^\alpha,v_{3,n_3}^\gamma)
       \hat\rho_f^{(\eta')}(v_{1,n_1}^{p_1},v_{2,n'}^\alpha,v_{3,n_3}^\gamma)}
 {\|v_{1,n_1}^{p_1}\|^2\|v_{2,n}^\alpha\|^2
  \|v_{2,n'}^\alpha\|^2\|v_{3,n_3}^\gamma\|^2}
 \nonumber\\
 &\quad\times\langle v_{2,n'}^\alpha|\Theta V[v_{1/2}(Q/2)](1)
       |v_{2,n}^\alpha\rangle
 \prod_{i=1}^2 F(h_{1,n_1}^{(i)},h_{2,n}^{(i)},h_{2,n'}^{(i)},h_{3,n_3}^{(i)};
 h_\psi^{(i)};c^{(i)}\mid q_1,\hat q_{2,1},\hat q_{2,2},q_3).
 \label{full-branching}
\end{align}
The additional split cut supplies one additional squared primary norm.
The quadratic part is the original theta sign; $(-1)^{2n_1}$ is the
enlarged NS contour sign.  The inserted tube operator is even, and both
split cuts carry the same total parity.  Using squared
norms in \eqref{full-branching} avoids a choice of complex square roots
in the numerical contraction.

For comparison with the normalized branching notation, define the
radially ordered coefficient, with the NS field in the middle slot, by
\ie
 \mathbb B(n',\tfrac12,n;\alpha,1-\alpha)
 =\frac{\langle v_{2,n'}^\alpha|V[v_{1/2}(Q/2)](1)
 |v_{2,n}^{1-\alpha}\rangle}
 {\|v_{2,n'}^\alpha\|\|v_{1/2}(Q/2)\|
  \|v_{2,n}^{1-\alpha}\|}.
 \label{radial-B}
\fe
No permutation of punctures is implicit in this notation.  The two
outer normalized coefficients use the original $(\NS,\R,\R)$ order.
For $M=2|n|-\frac12$, the raw primary definitions give
\ie
 \Theta v_n^0=-2^{(-1)^M/2}v_n^1,\qquad
 \Theta v_n^1=2^{-(-1)^M/2}v_n^0.
 \label{theta-primary}
\fe
With compatible norm roots such that
$\Theta(v_n^\alpha/\|v_n^\alpha\|)=-i
v_n^{1-\alpha}/\|v_n^{1-\alpha}\|$, anticommuting $\Theta$ through
$V[v_{1/2}]$ gives
\ie
 \frac{\langle v_{n'}^\alpha|\Theta V[v_{1/2}](1)|v_n^\alpha\rangle}
 {\|v_{n'}^\alpha\|\|v_n^\alpha\|}
 =i\|v_{1/2}\|\mathbb B(n',\tfrac12,n;\alpha,1-\alpha).
 \label{normalized-middle}
\fe
Accordingly the primary factor in \eqref{full-branching} is exactly
\ie
 &i\|v_{1/2}(Q/2)\|
 \mathbb B_f^{(\eta);p_1}(n_1,n,n_3;\alpha,\gamma)
 \mathbb B_f^{(\eta');p_1}(n_1,n',n_3;\alpha,\gamma)
 \mathbb B(n',\tfrac12,n;\alpha,1-\alpha).
\fe
Here $\|v_{1/2}(Q/2)\|^2=-Q^2$.  The preceding raw formula is the
definition when another square-root convention is used.

\section{The middle coefficient from the physical Ward identity}

The physical stress tensor acts on the SCA factor, whereas
\eqref{ordered-insertion} acts on the auxiliary factor.  Hence
\ie
 {}[L_m,\Theta Q\psi(1)]=0.
\fe
This statement concerns the physical $L_m$, not $L_m^{(1)}$ or
$L_m^{(2)}$.  BPZ conjugation, which is bilinear here, gives
\ie
 \langle L_1v_{n'}^\alpha|\Theta Q\psi(1)|v_n^\alpha\rangle
 &=\langle v_{n'}^\alpha|\Theta Q\psi(1)|L_{-1}v_n^\alpha\rangle,\\
 \langle L_{-1}v_{n'}^\alpha|\Theta Q\psi(1)|v_n^\alpha\rangle
 &=\langle v_{n'}^\alpha|\Theta Q\psi(1)|L_1v_n^\alpha\rangle.
 \label{middle-ward}
\fe
Use the reusable $\mathbb V$ action coefficients from the human notes:
\ie
 L_1v_n^\alpha&=\sum_{|A|+|B|=4n-3}\mathbb V_{1,n,\alpha}^{AB}
 L_{-A}^{(1)}L_{-B}^{(2)}v_{n-1}^\alpha,\\
 L_{-1}v_n^\alpha&=\sum_{i=1}^2\mathbb V_{n,\alpha}^{(i)}
 L_{-1}^{(i)}v_n^\alpha+
 \sum_{|A|+|B|=4n-1}\mathbb V_{-1,n,\alpha}^{AB}
 L_{-A}^{(1)}L_{-B}^{(2)}v_{n-1}^\alpha
 \qquad(n\geq3/4).
\fe
The levels follow from
$2n^2-2(n-1)^2=4n-2$, followed by the physical level change $\mp1$.
At the boundary $n=1/4$, $L_1v_{1/4}^\alpha=0$ and the second line is
replaced by its two same-primary level-one terms and a primary term
proportional to $v_{-3/4}^\alpha$.  Its coefficient is immaterial below
because that term is forbidden by fusion.  The action coefficients are
inputs from the finite double-Virasoro span decomposition in the human
notes; no independent recursion for those coefficients is assumed.
When $n'=n+\frac12>0$ and $n\geq\frac14$, $L_1$ on the larger
endpoint produces descendants of $v_{n'-1}^\alpha$.  In $L_{-1}v_n$
the other-primary branch has label $n-1$ (or $-3/4$ at $n=1/4$).
Its separation from $n'$ is $3/2$, so the middle fusion rule
\eqref{middle-fusion} kills its contribution.  The surviving same-primary
descendants obey
\ie
 \rho(\nu_{h'},\nu_{h_\psi},L_{-1}\nu_h)=h+h_\psi-h'.
\fe
The first Ward identity thus becomes the scalar recurrence
\begin{align}
 &\left[\sum_{i=1}^2\mathbb V_{n,\alpha}^{(i)}
 (h_n^{(i)}+h_\psi^{(i)}-h_{n'}^{(i)})\right]
 \langle v_{n'}^\alpha|\Theta Q\psi(1)|v_n^\alpha\rangle
 \nonumber\\
 &\quad=\sum_{|A|+|B|=4n'-3}\mathbb V_{1,n',\alpha}^{AB}
 \rho(L_{-A}\nu_{h_{n'-1}^{(1)}},\nu_{h_\psi^{(1)}},\nu_{h_n^{(1)}})
 \rho(L_{-B}\nu_{h_{n'-1}^{(2)}},\nu_{h_\psi^{(2)}},\nu_{h_n^{(2)}})
 \langle v_{n'-1}^\alpha|\Theta Q\psi(1)|v_n^\alpha\rangle.
 \label{middle-recurrence-out}
\end{align}
If $n=n'+\frac12>0$ and $n'\geq\frac14$, apply the second identity
instead.  Using
$\rho(L_{-1}\nu_{h'},\nu_{h_\psi},\nu_h)=h'+h_\psi-h$ gives
\begin{align}
 &\left[\sum_{i=1}^2\mathbb V_{n',\alpha}^{(i)}
 (h_{n'}^{(i)}+h_\psi^{(i)}-h_n^{(i)})\right]
 \langle v_{n'}^\alpha|\Theta Q\psi(1)|v_n^\alpha\rangle
 \nonumber\\
 &\quad=\sum_{|A|+|B|=4n-3}\mathbb V_{1,n,\alpha}^{AB}
 \rho(\nu_{h_{n'}^{(1)}},\nu_{h_\psi^{(1)}},L_{-A}\nu_{h_{n-1}^{(1)}})
 \rho(\nu_{h_{n'}^{(2)}},\nu_{h_\psi^{(2)}},L_{-B}\nu_{h_{n-1}^{(2)}})
 \langle v_{n'}^\alpha|\Theta Q\psi(1)|v_{n-1}^\alpha\rangle.
 \label{middle-recurrence-in}
\end{align}
For negative labels use the human notes' simultaneous reflection
$n\mapsto-n$, $P\mapsto-P$ in the action coefficients.  Applying $L_1$
to the endpoint of larger absolute label lowers that absolute label,
including the crossing $3/4\mapsto-1/4$.  Every allowed pair reaches
one of the two oriented ground pairs.  Their values, obtained from
\eqref{zero-mode-action}, follow directly from the ground primaries
\ie
 v_{\pm1/4}^0&=u^0\otimes w^+
 \mp e^{i\pi/4}u^1\otimes w^-,\\
 v_{\pm1/4}^1&=\frac1{\sqrt2}
 (u^1\otimes w^+\pm e^{i\pi/4}u^0\otimes w^-).
\fe
With the physical ground metric $\operatorname{diag}(1,i)$ and the
auxiliary metric $\operatorname{diag}(1,-1)$, their squared norms are
$2$ and $-1$.  Moreover $Q\Theta\psi_0$ maps $v_{\pm1/4}^0$ to
$(Q/\sqrt2)v_{\mp1/4}^0$, and $v_{\pm1/4}^1$ to
$-(Q/\sqrt2)v_{\mp1/4}^1$.  Therefore
\ie
 \langle v_{1/4}^0|\Theta Q\psi(1)|v_{-1/4}^0\rangle
 &=\langle v_{-1/4}^0|\Theta Q\psi(1)|v_{1/4}^0\rangle=\sqrt2Q,\\
 \langle v_{1/4}^1|\Theta Q\psi(1)|v_{-1/4}^1\rangle
 &=\langle v_{-1/4}^1|\Theta Q\psi(1)|v_{1/4}^1\rangle=Q/\sqrt2.
\fe
No other primary middle matrix elements are needed as anchors.
Equations \eqref{theta-primary} and \eqref{radial-B} then give
$\mathbb B(n',1/2,n;\alpha,1-\alpha)$.

At a generic momentum the scalar pivot in the recurrence is nonzero.
Exceptional zero pivots require continuation of the combined meromorphic
expression, not division by zero.  The implementation raises an error
at such a pivot.  Computing the finite ordinary three-point descendant
factors in these equations by Ward identities is distinct from computing
an entire conformal block.  All conformal blocks in
\eqref{full-branching} use the $c$-recursion described below.

\section{The convolution with the physical block}

Return temporarily to the tensor-product basis on the split tube.  Its
physical factor propagates through the identity.  A physical BPZ metric
at the middle sphere therefore cancels one of the two physical inverse
metrics.  Within each physical level and parity this is explicitly
$G^{-1}GG^{-1}=G^{-1}$.  The remaining physical state has equal level on
the two halves and propagation weight
\ie
 \hat q_{2,1}^{|B|}\hat q_{2,2}^{|B|}
 =q_2^{|B|}.
\fe
This contracts the physical graph back to the original theta graph.
The two contour phases in \eqref{outer-frame} multiply to $(-1)^A$:
the NS descendant parity and the total second-slot physical parity agree
at the paired vertices, so the two $f$-dependent signs cancel.  This
removes the physical part of the enlarged factor $(-1)^{A+\mathsf A}$.
The remaining $(-1)^{\mathsf A}$ is precisely the auxiliary NS contour
sign in the fermion block.  After this linear sign is accounted for,
only the quadratic theta signs remain to be compared.
The auxiliary factors retain the inserted field and independent levels
on the two halves.  Because $\Theta\psi$ is even, it identifies the
auxiliary parity bits on the two halves.  There is consequently no
fourth independent parity bit.  If $\epsilon_i$ and $\delta_i$ are the
physical and auxiliary parities, respectively, the theta sign at total
parity differs from the separate signs by
\ie
 &\sum_{i<j}(\epsilon_i+\delta_i)(\epsilon_j+\delta_j)
 -\sum_{i<j}\epsilon_i\epsilon_j-\sum_{i<j}\delta_i\delta_j\\
 &\qquad=\sum_{i<j}(\epsilon_i\delta_j+\delta_i\epsilon_j)
 \pmod2.
\fe
This gives the original Ramond convolution cocycle:
\ie
 \eta_1^{\epsilon_1}\eta_2^{\epsilon_2}\eta_3^{\epsilon_3}
 \star_{\R}\eta_1^{\delta_1}\eta_2^{\delta_2}\eta_3^{\delta_3}
 =(-1)^{\sum_{i<j}(\epsilon_i\delta_j+\delta_i\epsilon_j)}
 \eta_1^{\epsilon_1+\delta_1}\eta_2^{\epsilon_2+\delta_2}
 \eta_3^{\epsilon_3+\delta_3}.
 \label{cocycle}
\fe
Parity exponents are reduced modulo two.  The powers of sewing
parameters multiply ordinarily.  In particular the complete map on
monomials is
\begin{align}
 &(q_1^a q_2^b q_3^c\eta_1^{\epsilon_1}
       \eta_2^{\epsilon_2}\eta_3^{\epsilon_3})
 \star_{\R}
 (q_1^{a'}\hat q_{2,1}^{b'}\hat q_{2,2}^{b''}q_3^{c'}
       \eta_1^{\delta_1}\eta_2^{\delta_2}\eta_3^{\delta_3})
 \nonumber\\
 &\quad=(-1)^{\sum_{i<j}(\epsilon_i\delta_j+\delta_i\epsilon_j)}
 q_1^{a+a'}\hat q_{2,1}^{b+b'}\hat q_{2,2}^{b+b''}q_3^{c+c'}
 \eta_1^{\epsilon_1+\delta_1}\eta_2^{\epsilon_2+\delta_2}
 \eta_3^{\epsilon_3+\delta_3}.
 \label{coefficient-map}
\end{align}
Summing this identity over the physical and auxiliary descendants gives
\ie
 \widehat{\mathbb F}_f^{(\eta,\eta')}[\Theta v_{1/2}]
 =\mathbb F_f^{(\eta,\eta')}(q_1,\hat q_{2,1}\hat q_{2,2},q_3)
 \star_{\R}\mathbb F_{\mathsf F}[\Theta v_{1/2}]
 (q_1,\hat q_{2,1},\hat q_{2,2},q_3).
 \label{full-convolution}
\fe
This is an identity for the full inserted block.  Off-diagonal powers
of the two split parameters must be retained on both sides before
division.  Restricting them to equal powers would replace the inserted
field by its zero-mode contribution and would lose terms needed for
the convolution.

\subsection{The inverse in the required sector}

We first establish the sector identity from the physical sewing
definition.  On either Ramond module use the level-preserving replacement
\ie
 \mathbb L_{-B}w^+&\longmapsto
 (-1)^B e^{-i\pi/4}\mathbb L_{-B}w^-,\\
 \mathbb L_{-B}w^-&\longmapsto
 -(-1)^B e^{i\pi/4}\mathbb L_{-B}w^+.
\fe
Here $B$ in a sign counts the supercurrent generators, as in the human
notes.  The map is odd, supercommutes with the physical SCA, and
preserves the BPZ pairing with ground metric $\operatorname{diag}(1,i)$.
The last assertion follows at the ground level from
$(e^{-i\pi/4})^2 i=1$ and $(-e^{i\pi/4})^2=i$; contravariance extends
it to descendants.  The ground definitions of $\rho_f^{(\eta)}$ give
the factor $-i\eta(-1)^{\epsilon_2}$ when this replacement is applied
to both Ramond slots.  Including the Koszul sign for the two odd maps,
the simultaneous action intertwines the Ward identities.  The same
factor therefore holds for arbitrary descendants, where $\epsilon_2$
is the total physical parity of the second slot.

Applying this change of summation basis at both outer vertices multiplies
the unsigned sewing coefficient by $(-i\eta)(-i\eta')=-\eta\eta'$;
the two factors $(-1)^{\epsilon_2}$ cancel.  Flipping both Ramond parity
bits changes the quadratic theta sign by
$(-1)^{\epsilon_2+\epsilon_3+1}$.  Consequently, at each fixed triple
of levels, the coefficient at $(\epsilon_1,1-\epsilon_2,1-\epsilon_3)$
is $\eta\eta'(-1)^{\epsilon_2+\epsilon_3}$ times the coefficient at
$(\epsilon_1,\epsilon_2,\epsilon_3)$.  Equation \eqref{cocycle} also
gives
\ie
 (\eta_2\eta_3)\star_{\R}
 (\eta_1^{\epsilon_1}\eta_2^{\epsilon_2}\eta_3^{\epsilon_3})
 =(-1)^{\epsilon_2+\epsilon_3}
 \eta_1^{\epsilon_1}\eta_2^{1-\epsilon_2}\eta_3^{1-\epsilon_3}.
\fe
Thus $(\eta_2\eta_3)\star_{\R}\mathbb F_f^{(\eta,\eta')}
=\eta\eta'\mathbb F_f^{(\eta,\eta')}$ to every level.  In particular,
the opposite outer sector satisfies
\ie
 (\eta_2\eta_3)\star_{\R}\mathbb F_f^{(\eta,-\eta)}
 =-\mathbb F_f^{(\eta,-\eta)}.
 \label{sector-identity}
\fe
At zero level, the outer ground forms
$\rho_{\mathsf F}(\mathbf1,u^0,u^0)=1$ and
$\rho_{\mathsf F}(\mathbf1,u^1,u^1)=i$, together with
\eqref{zero-mode-action}, give the auxiliary constant
\ie
 \left.\mathbb F_{\mathsf F}[\Theta v_{1/2}]\right|_{q=0}
 =\frac{Q}{\sqrt2}(1-\eta_2\eta_3).
\fe
This constant has no inverse in the entire parity algebra.  However,
$(\eta_2\eta_3)\star_{\R}(\eta_2\eta_3)=1$, so its multiplication
on the sector \eqref{sector-identity} is the scalar $\sqrt2Q$.
For $Q\ne0$ this scalar is invertible.  Thus the original formal inverse
is replaced by a sector inverse whose identity is
$(1-\eta_2\eta_3)/2$:
\ie
 \mathbb F_{\mathsf F}^{-1}[\Theta v_{1/2}]
 \star_{\R}\mathbb F_{\mathsf F}[\Theta v_{1/2}]
 =\frac{1-\eta_2\eta_3}{2}.
\fe
Within this sector the physical block is uniquely determined by
\ie
 \mathbb F_f^{(\eta,-\eta)}
 =\mathbb F_{\mathsf F}^{-1}[\Theta v_{1/2}]
 \star_{\R}\widehat{\mathbb F}_f^{(\eta,-\eta)}[\Theta v_{1/2}].
 \label{recovery}
\fe
The computation need not construct an inverse as a separate series.
At each monomial subtract all convolutions involving already determined
lower coefficients and positive-level auxiliary coefficients; divide
the remainder by $\sqrt2Q$.  Equation \eqref{sector-identity} makes this
triangular division unique.  A remainder outside that sector is evidence
of incorrect input or insufficient numerical precision and must not be
silently projected away.

% Separate subsection: direct-definition auxiliary benchmark.
% The user selected this direct auxiliary backend after its benchmark.
\subsection{Direct fermion sewing: definition and level-ten benchmark}
\label{direct-fermion-definition}

The selected algorithm evaluates the auxiliary factor directly in its
fermion basis, while retaining central-charge recursion for both
Virasoro factors in the enlarged block.  This subsection derives the
direct auxiliary construction and records its measured cost for the
inserted genus-two graph.  The timing below concerns the auxiliary
factor alone and performs no comparison with another block; validation
of the integrated physical calculation is reported separately.

On the NS edge use the states
$\boldsymbol\Psi_{-\mathsf A}\mathbf1$, where $\mathsf A$ is a
strictly descending list of positive half-integers.  On each Ramond
edge use $\boldsymbol\Psi_{-\mathsf B}u^{\mathsf b}$, where
$\mathsf B$ is a strictly descending list of positive integers and
$\mathsf b=0,1$.  The symbol $|\mathsf B|$ denotes the sum of the
modes, whereas $\#\mathsf B$ counts them.  We retain the conventions
\begin{equation}
 \{\psi_r,\psi_s\}=\delta_{r+s,0},\qquad
 \psi_0u^{\mathsf b}=\frac{u^{1-\mathsf b}}{\sqrt2},\qquad
 \psi_r^\dagger=-\psi_{-r},\qquad
 \langle u^{\mathsf b}|u^{\mathsf b}\rangle=(-1)^{\mathsf b}.
\end{equation}
The Fock metric is diagonal.  Its diagonal entries are
$(-1)^{\#\mathsf A}$ in the NS sector and
$(-1)^{\#\mathsf B+\mathsf b}$ in the Ramond sector.

The incoming state on the additional sphere is at $z_3=0$ and carries
$\hat q_{2,1}^{|\mathsf B|}$; the outgoing state is at infinity and
carries $\hat q_{2,2}^{|\mathsf B'|}$.  Write their common parity and
the parities on the other two edges as
\begin{equation}
 \delta_1=\#\mathsf A,\qquad
 \delta_2=\#\mathsf B+\mathsf b
          =\#\mathsf B'+\mathsf b',\qquad
 \delta_3=\#\mathsf C+\mathsf c\pmod2.
\end{equation}
Both outer forms vanish unless
$\delta_1+\delta_2+\delta_3=0\pmod2$.

\paragraph{The complete middle insertion.}
At $z_3=1$ every term in
$\psi(z_3)=\sum_{m\in\mathbb Z}\psi_mz_3^{-m-1/2}$ has unit
coordinate factor.  Using
\begin{equation}
 \Theta\boldsymbol\Psi_{-\mathsf B}u^{\mathsf b}
 =(-1)^{\#\mathsf B+\mathsf b+1}
  \boldsymbol\Psi_{-\mathsf B}u^{1-\mathsf b},
\end{equation}
the zero-mode action is
\begin{equation}
 Q\Theta\psi_0\boldsymbol\Psi_{-\mathsf B}u^{\mathsf b}
 =\frac Q{\sqrt2}(-1)^{\mathsf b}
  \boldsymbol\Psi_{-\mathsf B}u^{\mathsf b}.
\end{equation}
For $k>0$, creation by $\psi_{-k}$, when $k\notin\mathsf B$, and
annihilation by $\psi_k$, when $k\in\mathsf B$, both give the sign
$(-1)^{\#\{m\in\mathsf B:m>k\}}$ in the descending basis.
The resulting word has one more or one fewer fermion.  The subsequent
$\Theta$ action therefore contributes $(-1)^{\delta_2}$ and changes
the ground label to $1-\mathsf b$.  Consequently the action
coefficients of the full insertion are
\begin{equation}
 [Q\Theta\psi(1)]_{(\mathsf B',\mathsf b'),(\mathsf B,\mathsf b)}
 =\begin{cases}
  \dfrac Q{\sqrt2}(-1)^{\mathsf b},
       &\mathsf B'=\mathsf B,\quad\mathsf b'=\mathsf b,\\[3pt]
  Q(-1)^{\#\{m\in\mathsf B:m>k\}+\delta_2},
       &\mathsf B'=\mathsf B\mathbin\triangle\{k\},
          \quad\mathsf b'=1-\mathsf b,\quad k>0,\\[3pt]
  0,&\text{otherwise}.
 \end{cases}
 \label{direct-fermion-middle}
\end{equation}
Here $\triangle\{k\}$ toggles the occupancy of $k$.  For each incoming
state, each nonzero mode connects to at most one outgoing Fock state.
Both transitions
preserve $\delta_2$, although their two propagation levels need not
agree.  Reversing a toggle leaves its sign unchanged: the number of
occupied modes larger than $k$ and the total parity are unchanged.

\paragraph{The state sum and its signs.}
In the following formula the subscript on
$[Q\Theta\psi(1)]$ means an action coefficient, as in
\eqref{direct-fermion-middle}, rather than a BPZ matrix element.
The latter contains the additional outgoing norm
$(-1)^{\delta_2}$.  The four inverse Fock norms multiply to
$(-1)^{\delta_1+2\delta_2+\delta_3}$, so including this outgoing
norm leaves $(-1)^{\delta_1+\delta_2+\delta_3}=1$ on every allowed
summand.  The NS contour sign and the original theta sign remain.
With the two outer forms in the order $(\NS,\R,\R)$, the complete
direct definition is therefore
\begin{align}
 &\mathbb F_{\mathsf F}[\Theta v_{1/2}]
 (q_1,\hat q_{2,1},\hat q_{2,2},q_3)\nonumber\\
 &=\sum_{\mathsf A,\mathsf B,\mathsf b,
                \mathsf B',\mathsf b',\mathsf C,\mathsf c}
 q_1^{|\mathsf A|}\hat q_{2,1}^{|\mathsf B|}
 \hat q_{2,2}^{|\mathsf B'|}q_3^{|\mathsf C|}
 \eta_1^{\delta_1}\eta_2^{\delta_2}\eta_3^{\delta_3}
 (-1)^{\delta_1+\delta_1\delta_2+\delta_1\delta_3+\delta_2\delta_3}
 \nonumber\\
 &\quad\times
 \rho_{\mathsf F}\!\left(
  \boldsymbol\Psi_{-\mathsf A}\mathbf1,
  \boldsymbol\Psi_{-\mathsf B}u^{\mathsf b},
  \boldsymbol\Psi_{-\mathsf C}u^{\mathsf c}\right)
 \rho_{\mathsf F}\!\left(
  \boldsymbol\Psi_{-\mathsf A}\mathbf1,
  \boldsymbol\Psi_{-\mathsf B'}u^{\mathsf b'},
  \boldsymbol\Psi_{-\mathsf C}u^{\mathsf c}\right)
 [Q\Theta\psi(1)]_{(\mathsf B',\mathsf b'),(\mathsf B,\mathsf b)}.
 \label{direct-fermion-sum}
\end{align}
The ground forms are
$\rho_{\mathsf F}(\mathbf1,u^0,u^0)=1$ and
$\rho_{\mathsf F}(\mathbf1,u^1,u^1)=i$.  They give the constant
$Q(1-\eta_2\eta_3)/\sqrt2$ directly.  In particular the second
split edge does not introduce an independent parity sign.

The outer forms are computed by moving one fermion contour at a time
between the punctures at infinity, one and zero.  For the three possible
target slots, the useful test sections are
\begin{equation}
 z^{k-1/2}(z-1)^{-1/2},\qquad
 z^{-1/2}(z-1)^{-n-1/2},\qquad
 z^{-n-1/2}(z-1)^{-1/2}.
\end{equation}
Their expansions give half-integer binomial coefficients and the
frame phases $(i,1,-i)$, together with the usual graded contour signs.
For the first section the target NS creation mode is $k-1/2$; in the
other two sections the target Ramond creation mode is $n$.  Remove
the largest occupied mode in the first nonempty slot and solve the
Ward relation for the original form.  Every remaining term has lower
total nonzero-mode level, so this procedure terminates at the displayed
ground forms.  Each local expansion is stopped only after its modes
annihilate every state on that leg.  The resulting bound is determined
by the occupied modes, rather than a fixed expansion order.
Computed three-point forms are stored and reused throughout the sum.

\paragraph{Exact rational arithmetic.}
No irrational arithmetic is needed after extracting $Q/\sqrt2$ from
the final block.  To see this, temporarily use $\sqrt2u^1$ in place
of $u^1$.  In this basis the two zero-mode coefficients are $1/2$
and $1$.  Strip the known phase $i^{\delta_2}$ from the outer form
and denote the stored value by $\rho_{\mathsf F}^{\mathrm{rat}}$.
Thus, with the same fermion words as above,
\begin{equation}
 \rho_{\mathsf F}
 =i^{\delta_2}2^{-(\mathsf b+\mathsf c)/2}
   \rho_{\mathsf F}^{\mathrm{rat}}.
\end{equation}
The stored ground values are $1$ and $2$.  The mode Ward relations
then involve only rational binomial coefficients and signs, so every
stored value is rational.  The two outer phases in a sewn term
multiply to $(-1)^{\delta_2}$.

For a diagonal transition, combining these phases, the factor
$(-1)^{\mathsf b}$ from $\Theta\psi_0$, and the two ground rescalings
gives the following normalized summand, apart from the common
propagation monomial and sewing sign already displayed in
\eqref{direct-fermion-sum}:
\begin{equation}
 \frac{(-1)^{\#\mathsf B}}{2^{\mathsf b+\mathsf c}}
  \left(\rho_{\mathsf F}^{\mathrm{rat}}
   (\mathsf A,\mathsf B,\mathsf C)\right)^2.
\end{equation}
For a toggle, $\mathsf b+\mathsf b'=1$.  The square root of two in
$[Q\Theta\psi(1)]/(Q/\sqrt2)$ cancels the remaining half-integer
ground-rescaling power, while the two factors
$(-1)^{\delta_2}$ cancel.  Its normalized summand is
\begin{equation}
 \frac{(-1)^{\#\{m\in\mathsf B:m>k\}}}{2^{\mathsf c}}
 \rho_{\mathsf F}^{\mathrm{rat}}(\mathsf A,\mathsf B,\mathsf C)
 \rho_{\mathsf F}^{\mathrm{rat}}(\mathsf A,\mathsf B',\mathsf C).
\end{equation}
The ground labels in these abbreviated arguments remain those in
\eqref{direct-fermion-sum}.  The temporary basis rescaling has
therefore canceled from the physical normalization.  The implementation
uses exact FLINT rational arithmetic and restores $Q/\sqrt2$ only
when a numerical value of the full block is requested.

\paragraph{Finite truncation and measured cost.}
For the balanced split
$\hat q_{2,1}=u\sqrt{q_2}$,
$\hat q_{2,2}=u^{-1}\sqrt{q_2}$, total level $N$ means
\begin{equation}
 |\mathsf A|+
 \frac{|\mathsf B|+|\mathsf B'|}{2}+|\mathsf C|\leq N.
\end{equation}
The stored exponent tuple is
$(2|\mathsf A|,|\mathsf B|,|\mathsf B'|,2|\mathsf C|)$;
its entries sum to at most $2N$.  Unequal split levels retain the
power $u^{|\mathsf B|-|\mathsf B'|}$.  At $N=10$ a single split
edge can reach level twenty.  To enumerate each nonzero-mode pair
once, start from its lower-level state and add the unique toggled
mode $k$.  The reverse removal supplies the transposed pair of
propagation powers with the same coefficient.  The diagonal part
is added separately.

A fresh-process run on September 10, 2026, using Python 3.14.3 on
macOS arm64, computed the complete normalized auxiliary series through
$N=10$ in $0.186281209\,\mathrm{s}$.  This construction time includes
the provider, Fock bases, Ward cache, all mode transitions, rational
normalization and coefficient accumulation.  Elapsed process time
through imports and the initial result save was
$0.290214667\,\mathrm{s}$; peak resident memory was
$46.765625\,\mathrm{MiB}$.  The saved result contains $4793$
nonzero exponent tuples and $9586$ nonzero rational parity
coefficients.  The calculation stored $7508$ three-point forms and
summed $20156$ nonzero state contractions.  Its coefficients and
timing metadata are recorded in
\path{Code/theta_fermion_ccy/results/direct_fermion_level10.json}.
The final small rewrite adding timing metadata is excluded from the
process-through-save measurement.  These measurements establish the
cost of this direct auxiliary calculation; they do not constitute a
numerical validation or a timing of the complete physical block.

Equation \eqref{direct-fermion-sum} gives the full auxiliary factor.
Its normalization relative to the reduced central-charge recursion is
explained in \S\ref{ccy-and-truncation}.

\paragraph{Other genera and plumbing channels.}
The ingredients in \eqref{direct-fermion-sum} are local: a fermion
Fock basis on each edge, its diagonal pairing, a three-point form on
each sphere, and the sparse action of any inserted fermion.  They can
be sewn on any pants graph with compatible NS and Ramond labels.
NS--NS--NS vertices use the corresponding fermion Ward forms;
NS--R--R vertices use the forms described above.  Spin transports
and graded reordering signs are taken from the chosen graph and
local coordinates.  The particular quadratic theta sign in
\eqref{direct-fermion-sum} belongs to the present channel and must
not be imposed unchanged on another graph.  A finite plumbing-level
cutoff again leaves a finite sum, and the local three-point data can
be reused.  This is an arbitrary-genus sewing construction, although
the current assembler and the reported runtime concern only the
specified four-edge genus-two graph.  Its computational cost can
grow with the genus and the graph.  This direct auxiliary evaluation
does not use an Ising Verma-module limit or an Ising central-charge
recursion.  The two ordinary Virasoro blocks in the enlarged numerator
continue to use the prescribed central-charge recursion.


% Component derivation for inclusion in the separate machine notes.
% This file is intentionally not a replacement for the manuscript.
\subsection{The Ising decomposition}

We keep the auxiliary Ramond parity on each cut.  The two irreducible
Virasoro modules in the Neveu--Schwarz Fock space have highest weights
$0$ and $1/2$.  Each Ramond parity gives an irreducible module of highest
weight $1/16$.  The field inserted on the additional sphere has weight
$1/2$.  The central charge in this subsection is that of the Ising
Virasoro algebra, namely $c=1/2$; it is not the superconformal central
charge.

The auxiliary three-point forms in the spin frames of the human notes
obey
\begin{align}
 \rho_{\mathsf F}(\mathbf1,u^0,u^0)&=1,&
 \rho_{\mathsf F}(\mathbf1,u^1,u^1)&=i,\\
 \rho_{\mathsf F}(\psi_{-1/2}\mathbf1,u^0,u^1)&=-\frac1{\sqrt2},&
 \rho_{\mathsf F}(\psi_{-1/2}\mathbf1,u^1,u^0)&=\frac{i}{\sqrt2}.
\end{align}
All other ground components vanish by parity.  To obtain the second
line, move the fermion contour at infinity onto the Ramond puncture at
one.  The Ward relation gives $i\rho_{\mathsf F}
(\psi_{-1/2}\mathbf1,u^{\mathsf b},u^{1-\mathsf b})+
\rho_{\mathsf F}(\mathbf1,\psi_0u^{\mathsf b},u^{1-\mathsf b})=0$.
Substitution of $\psi_0u^{\mathsf b}=u^{1-\mathsf b}/\sqrt2$ gives
the displayed values.

The ordered middle insertion acts on a ground state as
\begin{equation}
 Q\Theta\psi_0u^{\mathsf b}
 =\frac{Q}{\sqrt2}(-1)^{\mathsf b}u^{\mathsf b}.
\end{equation}
It commutes with parity and is a primary intertwiner of Ising weight
$1/2$, since $\Theta$ commutes with every auxiliary Virasoro generator.
Consequently its descendant matrix elements are precisely those in the
Virasoro one-point block.  Combining the two outer ground forms with
this middle matrix element and the theta sewing signs gives
the following sign ledger.  The middle raw pairing is $Q/\sqrt2$
in each row; the signs from its outgoing ground norm have already
been included.  The norm column is the product of the four internal
ground norms, and the last sign includes the NS fermion and theta
reordering signs:
\[
\begin{array}{c|cc|ccc}
 h_{\NS}&\mathsf b&\mathsf c&\rho_{\mathsf F}^{2}
 &\text{norm product}&\text{sewing sign}\\\hline
 0&0&0&1&1&1\\
 0&1&1&-1&-1&-1\\
 \tfrac12&0&1&\tfrac12&1&1\\
 \tfrac12&1&0&-\tfrac12&-1&1
\end{array}
\]
Thus the complete Ising decomposition is
\begin{align}
 \mathbb F_{\mathsf F}[\Theta v_{1/2}]
 ={}&\frac{Q}{\sqrt2}(1-\eta_2\eta_3)
 F\left(0,\frac1{16},\frac1{16},\frac1{16};\frac12;\frac12
 \mathrel\big|q_1,\hat q_{2,1},\hat q_{2,2},q_3\right)\nonumber\\
 &+\frac{Q}{2\sqrt2}(\eta_1\eta_2+\eta_1\eta_3)q_1^{1/2}
 F\left(\frac12,\frac1{16},\frac1{16},\frac1{16};\frac12;\frac12
 \mathrel\big|q_1,\hat q_{2,1},\hat q_{2,2},q_3\right).
 \label{eq:ising-ordered-theta-decomposition}
\end{align}
Both blocks on the right must use the \emph{irreducible} Ising modules.
Their primary coefficients are one.  The four internal entries refer,
in order, to the Neveu--Schwarz edge, the two parts of the split Ramond
edge, and the remaining Ramond edge.  The next entry is the external
weight and the last is the central charge.  The $q_1^{1/2}$ in the
second term restores the physical level of the Neveu--Schwarz fermion.
This formula contains the complete parity dependence; all remaining
coefficients are ordinary scalar Virasoro coefficients.


\section{Central-charge recursion and truncation}
\label{ccy-and-truncation}

For generic internal weights, \eqref{punctured-vir} is computed with the
central-charge recursion of Cho, Collier and Yin~\cite{CCY}.  Its poles
come from a null state on one of its four internal edges.  At a pole on
edge $e$, the internal weight shifts by $rs$, the corresponding sewing
factor contributes $q_e^{rs}$, and the residue multiplies the lower
block by
\ie
 \frac{-\partial_{h_e}c_{rs}(h_e) A_{rs}}
 {c-c_{rs}(h_e)}
 P_{rs}[\text{spectators at the first endpoint}]
 P_{rs}[\text{spectators at the second endpoint}].
\fe
For the ordered graph \eqref{punctured-vir}, the two spectator pairs are
\[
\begin{array}{c|cc}
e&\text{first pair}&\text{second pair}\\\hline
1&(h_3,h_L)&(h_3,h_R)\\
L&(h_3,h_1)&(h_R,h_\psi)\\
R&(h_3,h_1)&(h_L,h_\psi)\\
3&(h_1,h_L)&(h_1,h_R)
\end{array}
\]
The order in each pair matters for odd null level.  The external
$h_\psi$ stays fixed during the recursion, including when its initial
value is degenerate.  Replacing it by its degenerate formula evaluated
at the recursive central charge would compute a different function.

For completeness the pole and residue data are the ordinary Virasoro
ones used in the human notes:
\ie
 b_{rs}(h)^2&=\frac{rs-1+2h\pm
 \sqrt{(r-s)^2+4(rs-1)h+4h^2}}{1-r^2},\\
 c_{rs}(h)&=1+6(b_{rs}+b_{rs}^{-1})^2,\\
 A_{rs}&=\frac12\prod_{\substack{m=1-r,\ldots,r\,,\ k=1-s,\ldots,s\\
 (m,k)\ne(0,0),(r,s)}}(mb_{rs}+kb_{rs}^{-1})^{-1},\\
 P_{rs}(h_a,h_b)&=
 \prod_{\substack{j=1-r,\,3-r,\ldots,r-1\\ k=1-s,\,3-s,\ldots,s-1}}
 \frac{\lambda_a+\lambda_b+jb_{rs}+kb_{rs}^{-1}}2
 \frac{\lambda_a-\lambda_b+jb_{rs}+kb_{rs}^{-1}}2,
\fe
where $h_a=((b_{rs}+b_{rs}^{-1})^2-\lambda_a^2)/4$ and likewise for
$h_b$.  The sum runs over $r\geq2,s\geq1$.  The numerical continuation
chooses the noncancelling quadratic root (the larger absolute value),
consistently across $(r,s)$ and $(s,r)$.  This matters at negative weights
when $s=1$, where a cancelling root can be zero.  Unresolved confluent
poles raise an error; a singular residue is never simply discarded.

The regular term at large $c$ is the global $SL(2)$ block times the
universal large-$c$ vacuum factor.  The global coefficient at levels
$(a,b,c,d)$ is
\ie
 \frac{\rho(L_{-1}^a\nu_{h_1},L_{-1}^b\nu_{h_L},L_{-1}^d\nu_{h_3})
       \rho(L_{-1}^a\nu_{h_1},L_{-1}^c\nu_{h_R},L_{-1}^d\nu_{h_3})
       \rho(L_{-1}^c\nu_{h_R},\nu_{h_\psi},L_{-1}^b\nu_{h_L})}
 {a!(2h_1)_a\,b!(2h_L)_b\,c!(2h_R)_c\,d!(2h_3)_d}.
\fe
The universal factor is independent of all light weights.  We may
therefore forget the external puncture and collapse the extra sphere.
Every sewing and collapsing map is M\"obius, so no Schwarzian term is
introduced.  Equation \eqref{product} shows that the factor depends
only on $(q_1,\hat q_{2,1}\hat q_{2,2},q_3)$.
It is computed by the leading Gaussian contractions of
$L_{-m}/\sqrt c$, $m\geq2$.  For multiplicities $k_m$ the diagonal norm is
\ie
 \left\|\prod_{m\geq2}(L_{-m}/\sqrt c)^{k_m}\mathbf1\right\|^2
 \longrightarrow
 \prod_{m\geq2} k_m!\left[\frac{m(m^2-1)}{12}\right]^{k_m}.
\fe
The pair contractions between the ordered slots are
\ie
 C_{\infty,0}(m,n)&=\delta_{mn}\frac{m(m^2-1)}{12},\\
 C_{\infty,1}(m,n)&=\frac{m(m^2-1)}{12}\binom{m-2}{n-2}
 \quad(m\geq n),\\
 C_{1,0}(m,n)&=\frac{(-1)^m(m+n-1)!}{12(m-2)!(n-2)!}.
\fe
The second is zero when $m<n$, and same-slot contractions are zero.
These formulas follow by expanding or differentiating
$\langle T(z)T(w)\rangle=c/[2(z-w)^4]$ at the three punctures.
At a fixed triple of vacuum levels, sum the square of the three-slot
Wick contraction over oscillator partitions and divide by the three
diagonal norms above.  The code performs this finite sum exactly over
the rationals.  It uses the large-$c$ algebra, not a finite-$c$ Gram
inverse.  The reduced double-Virasoro numerator is missing two copies
of this factor.  The selected direct fermion sum contains its full
sewing dependence, so division by it must be followed by restoring
two copies in the physical result.

To define a finite physical total-level truncation without placing the
insertion at an endpoint, write
\ie
 \hat q_{2,1}=u\sqrt{q_2},\qquad
 \hat q_{2,2}=u^{-1}\sqrt{q_2}.
\fe
A term $q_1^{a/2}\hat q_{2,1}^l\hat q_{2,2}^r q_3^{d/2}$ has
physical total degree $(a+l+r+d)/2$ and insertion-position dependence
$u^{l-r}$.  Through physical total level $N$ retain all terms with
$a+l+r+d\leq2N$.  This is a downward-closed truncation, so the
triangular convolution uses no coefficients above the cutoff.
The physical answer must have $l=r$ after division, but this property
is a result to be checked, not a restriction imposed on the computation.
At $N=10$ the contributing middle primary pairs satisfy
$|n|,|n'|\leq9/4$.  For the middle recurrence the reusable $L_1$ and
relevant $L_{-1}$ coefficients require descendant levels at most six.
This bound does not apply to the entire outer Ward grid: the current
outer implementation also constructs $L_{-1}v_{9/4}$, whose other-primary
branch reaches descendant level eight.
For negative branch labels, the implementation first uses the
reflection $v_{-n}(P)\leftrightarrow v_n(-P)$ and reverses the labels
of every term in the action expansion.  This is an intertwiner of
the physical SCA and both Virasoro algebras.  On the physical Ramond
grounds it acts as $\operatorname{diag}(1,-1)$ and maps the defining
fermion strings without an extra normalization factor.  Consequently
the action coefficients can be solved on the positive-label chart;
no high-level Fock-to-PBW reflection matrix is needed.  The action
data are stored for each momentum and reused across the outer Ward
equations.

\section{Implementation and numerical status}

The implementation is in \texttt{Code/theta\_fermion\_ccy}.
\texttt{middle\_branching.py} implements
\eqref{middle-recurrence-out}--\eqref{middle-recurrence-in};
\texttt{outer\_branching.py} uses the existing outer $L_1$ Ward system;
\texttt{punctured\_ccy.py} supplies genuine central-charge residues and
the global seed; \texttt{direct\_fermion.py} evaluates the full
auxiliary factor by its mode Ward identities and diagonal Fock
pairing; \texttt{pipeline.py} performs the branch sum and
\texttt{series\_algebra.py} performs the sector division.
The archived outer oracle uses the opposite label for an odd form:
$\eta_{\mathrm{native}}=(-1)^f\eta_{\mathrm{physical}}$.  The adapter
performs this transport before solving its Ward system.  It follows
algebraically from comparing the two $G_{-1/2}$ Ward conventions; the
numerical benchmark below tests $f=0$ only.

Only the two requested numerical comparisons are designated as
validation: the final physical block against an independent physical
PBW calculation through total level five, and the final physical block
at different split parameters with their product held fixed.  The PBW
reference is stored separately and is not imported by the production
pipeline.  No additional component numerical comparison suite is used.

\textbf{The two requested comparisons pass through level five.}
The benchmark is $b=7/5$, $(P_1,P_2,P_3)=(11/23,13/29,17/31)$,
$p_1=f=0$ and $(\eta,\eta')=(+,-)$.  The calculation uses 70 decimal
digits for the outer Ward system, 100 for the CCY engines, and exact
rational fermion sewing before restoring its common $Q/\sqrt2$ factor.  The
independent physical PBW reference uses 384-bit FLINT arithmetic.
All 4,648 parity-coefficient slots in the four-variable truncation are
compared, including off-diagonal slots whose physical answer is zero.
The maximum absolute difference is $4.72\times10^{-52}$; the maximum
difference divided by $\max(1,|\text{PBW coefficient}|)$ is
$4.72\times10^{-52}$.

The split-position comparison uses $(q_1,q_2,q_3)=(0.013,0.007,0.009)$,
$u=1/2,1,2$, and all eight evaluations of the spin-lift signs.  The
maximum change, divided by $\max(1,|\mathbb F(u=1)|)$, is
$2.96\times10^{-60}$.  For example, at $(\eta_1,\eta_2,\eta_3)=(1,-1,1)$
all three values agree with $2.0178804284272157747114188332429722451467$
to the digits shown.
These are numerical differences, not interval error bounds.

The pipeline timer includes all mathematical stages and intermediate
checkpoint writes.  It starts after imports and argument parsing and
ends before encoding and saving the final physical series.
This level-five run took $82.1968$ seconds in total: $52.1231$ for the
outer coefficients, $2.6851$ for the middle recurrence, $26.1996$ for
the two Virasoro blocks in the branch sum, $0.0825$ for the direct
auxiliary factor, and $0.6678$ for division and the restored vacuum
factors.  The balance is
assembly and intermediate serialization.  This is a \emph{level-five}
runtime; it is not a level-ten estimate.  The complete output and the
two comparisons are saved under \texttt{Code/theta\_fermion\_ccy/results}
as \texttt{ccy\_direct\_reflected\_level5.json} and
\texttt{validation\_ccy\_direct\_reflected\_level5.json}.

The production calculation retains the input, anchor and Ward-solution
precision throughout.  Its sector guard raises an error if division
leaves the required ideal; it does not project the answer to force the
identity.  The separate level-ten auxiliary benchmark in
\S\ref{direct-fermion-definition} is not a runtime for the physical
block.  The full level-ten production timing is being measured with
the same implementation and precision settings.

\begin{thebibliography}{9}
\bibitem{CCY}
M. Cho, S. Collier and X. Yin, \emph{Recursive Representations of
Arbitrary Virasoro Conformal Blocks},
\href{https://arxiv.org/abs/1703.09805}{arXiv:1703.09805}.
\bibitem{minimal-recursion}
N. Javerzat, R. Santachiara and O. Foda, \emph{Notes on the solutions
of Zamolodchikov-type recursion relations in Virasoro minimal models},
\href{https://arxiv.org/abs/1806.02790}{arXiv:1806.02790}.
\end{thebibliography}
\end{document}
